\begin{proof}[Proof of \cref{obj:thm:uniform-baseline}]
For admissible \(n,m,d,\epsilon\) and \(P\in\mathcal M_{d,\epsilon}\), define
\[
\rho_{\mathrm B}
=
\frac{1}{n\epsilon}
+
\left(\frac{d}{(n+m)\epsilon}\right)^2,
\qquad
\mathcal R_{\mathrm B}(P)
=
\mathbb E_P\!\left[
\bigl(T^{\mathrm B}(\mathcal D)-\tau(P)\bigr)^2
\right].
\]
The conditional outcome means belong to \([0,1]\), while the covariate masses are nonnegative and sum to one. Thus \cref{obj:def:target} gives
\[
-1
=
\sum_{j=1}^d p_j(P)(-1)
\le \tau(P)
\le
\sum_{j=1}^d p_j(P)
=
1.
\]
We first establish an uncapped risk bound for \(n\ge6\).

\begin{enumerate}
\item
Set, as in \cref{obj:def:baseline},
\[
h_{\mathrm o}=\lfloor n/2\rfloor,
\qquad
h_{\mathrm m}=n-h_{\mathrm o}+m,
\qquad
u_{\mathrm B}=h_{\mathrm o}/8,
\qquad
t_{\mathrm B}=h_{\mathrm m}/8.
\]
Both pool lengths and both prefix means are positive. By
\cref{obj:ass:data-product}, the outcome pool and the marginal pool are independent, with respective laws
\[
P^{\otimes h_{\mathrm o}}
\quad\text{and}\quad
P_{XA}^{\otimes h_{\mathrm m}}.
\]
Indeed, the marginal pool uses projections of complete records disjoint from the outcome pool, followed by the independent auxiliary records.

Extend these pools by independent iid continuations, and write
\[
\mathcal S_{\mathrm o}
=
\bigl((X_i^{\mathrm o},A_i^{\mathrm o},Y_i^{\mathrm o})\bigr)_{i\ge1}
\sim P^{\otimes\mathbb N},
\qquad
\mathcal S_{\mathrm m}
=
\bigl((X_i^{\mathrm m},A_i^{\mathrm m})\bigr)_{i\ge1}
\sim P_{XA}^{\otimes\mathbb N}.
\]
Independently of these sequences, take independent prefix lengths
\[
J_{\mathrm o}\sim\operatorname{Poi}(u_{\mathrm B}),
\qquad
J_{\mathrm m}\sim\operatorname{Poi}(t_{\mathrm B}),
\]
where \(\operatorname{Poi}(\lambda)\) is the Poisson probability law with mean \(\lambda\). Define the complete-atom and marginal-atom counts by
\[
N_{jay}^{\mathrm o}
=
\sum_{i=1}^{J_{\mathrm o}}
\mathbf1\{X_i^{\mathrm o}=j,A_i^{\mathrm o}=a,Y_i^{\mathrm o}=y\},
\qquad
N_{ja}^{\mathrm m}
=
\sum_{i=1}^{J_{\mathrm m}}
\mathbf1\{X_i^{\mathrm m}=j,A_i^{\mathrm m}=a\},
\]
for \(j\in\{1,\ldots,d\}\) and \(a,y\in\{0,1\}\); empty sums equal zero.

For a nonnegative integer array
\(\boldsymbol c=(c_{jay})\), put
\[
|\boldsymbol c|=\sum_{j,a,y}c_{jay}.
\]
Conditioning on the outcome-prefix length gives its Poisson probability multiplied by the multinomial probability:
\[
\begin{aligned}
\Pr\!\left((N_{jay}^{\mathrm o})=\boldsymbol c\right)
&=
e^{-u_{\mathrm B}}
\frac{u_{\mathrm B}^{|\boldsymbol c|}}{|\boldsymbol c|!}
\frac{|\boldsymbol c|!}{\prod_{j,a,y}c_{jay}!}
\prod_{j,a,y}
P(X=j,A=a,Y=y)^{c_{jay}}\\
&=
\prod_{j,a,y}
e^{-u_{\mathrm B}P(X=j,A=a,Y=y)}
\frac{
\bigl(u_{\mathrm B}P(X=j,A=a,Y=y)\bigr)^{c_{jay}}
}{c_{jay}!}.
\end{aligned}
\]
Here \(0^0=1\). The same calculation for a nonnegative integer array
\(\boldsymbol k=(k_{ja})\) gives
\[
\Pr\!\left((N_{ja}^{\mathrm m})=\boldsymbol k\right)
=
\prod_{j,a}
e^{-t_{\mathrm B}s_{aj}(P)}
\frac{\bigl(t_{\mathrm B}s_{aj}(P)\bigr)^{k_{ja}}}{k_{ja}!}.
\]
Consequently, the atom counts within each pool are independent Poisson variables, and independence of the pools gives independence across the two collections.

In particular, define
\[
Z_{aj}^{\mathrm B}=N_{ja1}^{\mathrm o},
\qquad
K_{aj}^{\mathrm B}=N_{ja}^{\mathrm m}.
\]
Then
\[
Z_{aj}^{\mathrm B}\sim
\operatorname{Poi}\!\bigl(u_{\mathrm B}q_{aj}(P)\bigr),
\qquad
K_{aj}^{\mathrm B}\sim
\operatorname{Poi}\!\bigl(t_{\mathrm B}s_{aj}(P)\bigr).
\]
For each fixed arm \(a\), the collection
\[
\bigl(Z_{aj}^{\mathrm B},K_{aj}^{\mathrm B},
K_{1-a,j}^{\mathrm B}\bigr)_{j=1}^d
\]
is mutually independent. These conclusions include zero-mass cells, whose counts equal zero almost surely.
% lean: CausalSmith.Stat.AnnotationRarearmFrontier.baseline_prefix_counts_law
% lean: CausalSmith.Stat.AnnotationRarearmFrontier.baseline_prefix_arm_counts_independent

\item
Define the arm totals, their population targets, and the clipped Poisson-prefix contrast by
\[
H_{\mathrm{sum},a}^{\mathrm B}
=
\sum_{j=1}^d
\frac{Z_{aj}^{\mathrm B}}{u_{\mathrm B}}
\left(1+\frac{K_{1-a,j}^{\mathrm B}}{K_{aj}^{\mathrm B}+1}\right),
\qquad
\psi_a
=
\sum_{j=1}^d p_j(P)\mu_{aj}(P),
\]
\[
T_{\mathrm{Pois}}^{\mathrm B}
=
\max\!\left\{-1,\,
\min\!\left\{1,\,
H_{\mathrm{sum},1}^{\mathrm B}
-
H_{\mathrm{sum},0}^{\mathrm B}
\right\}\right\}.
\]
The count laws and independence just established discharge the hypotheses of
\cref{obj:lem:inverse-count-arm-risk}. Let \(C_{\mathrm{arm}}>0\) be its universal variance constant. For each \(a\in\{0,1\}\), it supplies
\[
\left|
\mathbb E H_{\mathrm{sum},a}^{\mathrm B}-\psi_a
\right|
\le
\min\!\left\{
1,\frac{d}{\exp(1)t_{\mathrm B}\epsilon}
\right\},
\]
\[
\operatorname{Var}\!\left(H_{\mathrm{sum},a}^{\mathrm B}\right)
\le
C_{\mathrm{arm}}
\left(
\frac{1}{u_{\mathrm B}\epsilon}
+
\frac{1}{t_{\mathrm B}\epsilon}
\right).
\]
The finite arm sums are square-integrable, using Poisson moments and the positive denominators. Expanding squared loss around the mean therefore yields
\[
\begin{aligned}
\mathbb E\!\left[
\bigl(H_{\mathrm{sum},a}^{\mathrm B}-\psi_a\bigr)^2
\right]
&=
\operatorname{Var}\!\left(H_{\mathrm{sum},a}^{\mathrm B}\right)
+
\left(\mathbb E H_{\mathrm{sum},a}^{\mathrm B}-\psi_a\right)^2\\
&\le
C_{\mathrm{arm}}
\left(
\frac{1}{u_{\mathrm B}\epsilon}
+
\frac{1}{t_{\mathrm B}\epsilon}
\right)
+
\left(\frac{d}{\exp(1)t_{\mathrm B}\epsilon}\right)^2.
\end{aligned}
\]
By \cref{obj:def:target},
\[
\tau(P)=\psi_1-\psi_0.
\]
Projection onto \([-1,1]\) decreases squared distance to \(\tau(P)\). Moreover,
\[
\begin{aligned}
\bigl(T_{\mathrm{Pois}}^{\mathrm B}-\tau(P)\bigr)^2
&\le
\bigl(H_{\mathrm{sum},1}^{\mathrm B}
-H_{\mathrm{sum},0}^{\mathrm B}-\tau(P)\bigr)^2\\
&\le
2\bigl(H_{\mathrm{sum},1}^{\mathrm B}-\psi_1\bigr)^2
+
2\bigl(H_{\mathrm{sum},0}^{\mathrm B}-\psi_0\bigr)^2,
\end{aligned}
\]
where the second inequality follows by expanding the square and using the nonnegativity of
\[
\left[
\bigl(H_{\mathrm{sum},1}^{\mathrm B}-\psi_1\bigr)
+
\bigl(H_{\mathrm{sum},0}^{\mathrm B}-\psi_0\bigr)
\right]^2.
\]
Set
\[
C_{\mathrm{prefix}}=4C_{\mathrm{arm}}>0.
\]
Taking expectations gives
\[
\mathbb E\!\left[
\bigl(T_{\mathrm{Pois}}^{\mathrm B}-\tau(P)\bigr)^2
\right]
\le
C_{\mathrm{prefix}}
\left(
\frac{1}{u_{\mathrm B}\epsilon}
+
\frac{1}{t_{\mathrm B}\epsilon}
\right)
+
4\left(\frac{d}{\exp(1)t_{\mathrm B}\epsilon}\right)^2.
\]
% lean: CausalSmith.Stat.AnnotationRarearmFrontier.baseline_poisson_arm_mse
% lean: CausalSmith.Stat.AnnotationRarearmFrontier.baseline_poisson_contrast_risk
% lean: CausalSmith.Stat.AnnotationRarearmFrontier.baseline_ordered_prefix_risk

\item
Define the Poisson risk allowance and the exponential overflow allowance by
\[
\mathcal A_{\mathrm B}
=
C_{\mathrm{prefix}}
\left(
\frac{1}{u_{\mathrm B}\epsilon}
+
\frac{1}{t_{\mathrm B}\epsilon}
\right)
+
4\left(\frac{d}{\exp(1)t_{\mathrm B}\epsilon}\right)^2,
\qquad
\mathcal E_{\mathrm B}
=
e^{-h_{\mathrm o}}+e^{-h_{\mathrm m}}.
\]
Both are nonnegative. The clipped prefix statistic is measurable and takes values in \([-1,1]\), including on empty prefixes. The target also belongs to this interval. Thus the two positive, independent iid pools satisfy the hypotheses of
\cref{obj:lem:finite-prefix-transfer}.

More explicitly, writing \(\mathcal S^{1:h}\) for the first \(h\) records of a sequence, define
\[
T_{\mathrm{fin}}^{\mathrm B}
=
\mathbb E\!\left[
T_{\mathrm{Pois}}^{\mathrm B}
\mathbf1\{J_{\mathrm o}\le h_{\mathrm o},\,
J_{\mathrm m}\le h_{\mathrm m}\}
\,\middle|\,
\mathcal S_{\mathrm o}^{1:h_{\mathrm o}},
\mathcal S_{\mathrm m}^{1:h_{\mathrm m}}
\right].
\]
Expanding this conditional expectation over the independent Poisson lengths gives exactly the finite average in \cref{obj:def:baseline}, with zero output on overflow. Hence, after forming the pools from the original records,
\[
T^{\mathrm B}(\mathcal D)=T_{\mathrm{fin}}^{\mathrm B},
\qquad
\mathcal R_{\mathrm B}(P)
=
\mathbb E\!\left[
\bigl(T_{\mathrm{fin}}^{\mathrm B}-\tau(P)\bigr)^2
\right].
\]
Integration over the independent seed leaves this risk unchanged.

The transfer bound in \cref{obj:lem:finite-prefix-transfer}, with its nonnegative overflow allowance enlarged by a factor of four, now gives
\[
\begin{aligned}
\mathcal R_{\mathrm B}(P)
&\le
\mathbb E\!\left[
\bigl(T_{\mathrm{Pois}}^{\mathrm B}-\tau(P)\bigr)^2
\right]
+
4\mathcal E_{\mathrm B}\\
&\le
C_{\mathrm{prefix}}
\left(
\frac{1}{u_{\mathrm B}\epsilon}
+
\frac{1}{t_{\mathrm B}\epsilon}
\right)
+
4\left(\frac{d}{\exp(1)t_{\mathrm B}\epsilon}\right)^2
+
4\left(e^{-h_{\mathrm o}}+e^{-h_{\mathrm m}}\right)\\
&=
\mathcal A_{\mathrm B}+4\mathcal E_{\mathrm B}.
\end{aligned}
\]
% lean: CausalSmith.Stat.AnnotationRarearmFrontier.baseline_ruleRisk_eq_fixed_average
% lean: Causalean.Stat.FiniteRaoBlackwell.IndependentPoissonPrefix.sqRisk_independentPrefixRaoBlackwellStatistic_le

\item
The integer pool sizes satisfy
\[
n\le3\lfloor n/2\rfloor,
\qquad
n+m\le2\bigl(n-\lfloor n/2\rfloor+m\bigr),
\]
because \(n\ge6\) and \(m\ge0\). Consequently,
\[
h_{\mathrm o}\ge\frac n3,
\qquad
h_{\mathrm m}\ge\frac{n+m}{2},
\]
and therefore
\[
\frac{1}{u_{\mathrm B}\epsilon}
\le\frac{24}{n\epsilon},
\qquad
\frac{1}{t_{\mathrm B}\epsilon}
\le\frac{16}{n\epsilon}.
\]
Also, \(\exp(1)\ge1\) and \(t_{\mathrm B}\ge(n+m)/16\), so
\[
\frac{d}{\exp(1)t_{\mathrm B}\epsilon}
\le
16\frac{d}{(n+m)\epsilon}.
\]

For the overflow allowance, both pool sizes are at least \(n/3\), giving
\[
\mathcal E_{\mathrm B}\le2e^{-n/3}.
\]
The elementary bound \(x e^{-x}\le e^{-1}\), obtained by maximizing \(x e^{-x}\), together with \(\epsilon\le1/4\), yields
\[
(n\epsilon)\,2e^{-n/3}
\le
\frac n2e^{-n/3}
=
\frac32\left(\frac n3e^{-n/3}\right)
\le
\frac32e^{-1}
\le2.
\]
Since \(n\epsilon>0\), this proves
\[
\mathcal E_{\mathrm B}\le\frac{2}{n\epsilon}.
\]
Combining these estimates gives
\[
\begin{aligned}
\mathcal A_{\mathrm B}+\mathcal E_{\mathrm B}
&\le
\frac{40C_{\mathrm{prefix}}+2}{n\epsilon}
+
1024\left(\frac{d}{(n+m)\epsilon}\right)^2\\
&\le
(40C_{\mathrm{prefix}}+1026)\rho_{\mathrm B}.
\end{aligned}
\]
% lean: CausalSmith.Stat.AnnotationRarearmFrontier.baseline_pool_sizes
% lean: CausalSmith.Stat.AnnotationRarearmFrontier.baseline_pool_overflow_bound
% lean: CausalSmith.Stat.AnnotationRarearmFrontier.baseline_pool_rate_bound

\item
Define
\[
C_{\mathrm{large}}
=
4(40C_{\mathrm{prefix}}+1026)>0.
\]
Using \(\mathcal A_{\mathrm B}\ge0\), the preceding bounds imply, for every \(n\ge6\),
\[
\begin{aligned}
\mathcal R_{\mathrm B}(P)
&\le
\mathcal A_{\mathrm B}+4\mathcal E_{\mathrm B}\\
&\le
4(\mathcal A_{\mathrm B}+\mathcal E_{\mathrm B})\\
&\le
C_{\mathrm{large}}\rho_{\mathrm B}.
\end{aligned}
\]
The constants used here are universal.
% lean: CausalSmith.Stat.AnnotationRarearmFrontier.uniform_baseline

\item
Choose the final constant
\[
C=\max\{C_{\mathrm{large}},4\}>0.
\]
For all admissible sample sizes, the rule in \cref{obj:def:baseline} is measurable: the small-sample branch is constant, and the remaining branch is a finite sum of measurable count functions with positive denominators.

To check its range directly, for \(n\ge6\) define
\[
w_{\mathrm B}(h',h'')
=
e^{-u_{\mathrm B}-t_{\mathrm B}}
\frac{u_{\mathrm B}^{h'}}{h'!}
\frac{t_{\mathrm B}^{h''}}{h''!},
\qquad
0\le h'\le h_{\mathrm o},\quad
0\le h''\le h_{\mathrm m}.
\]
These weights are nonnegative and satisfy
\[
\sum_{h'=0}^{h_{\mathrm o}}
\sum_{h''=0}^{h_{\mathrm m}}
w_{\mathrm B}(h',h'')
=
\Pr(J_{\mathrm o}\le h_{\mathrm o})\,
\Pr(J_{\mathrm m}\le h_{\mathrm m})
\le1.
\]
Every clipped prefix value belongs to \([-1,1]\). Thus
\[
-1
\le
-\sum_{h',h''}w_{\mathrm B}(h',h'')
\le
T^{\mathrm B}(\mathcal D)
\le
\sum_{h',h''}w_{\mathrm B}(h',h'')
\le1.
\]
For \(n<6\), the rule equals zero, so the same range conclusion holds.
% lean: CausalSmith.Stat.AnnotationRarearmFrontier.baselineEstimator_measurable
% lean: CausalSmith.Stat.AnnotationRarearmFrontier.baselineEstimator_mem_Icc

\item
Suppose \(1\le n<6\). The zero rule has risk
\[
\mathcal R_{\mathrm B}(P)=\tau(P)^2\le1.
\]
Moreover,
\[
0<n\epsilon\le\frac54,
\qquad
\frac{1}{n\epsilon}\ge\frac45,
\qquad
\min\{1,\rho_{\mathrm B}\}\ge\frac45.
\]
It follows that
\[
\mathcal R_{\mathrm B}(P)
\le1
\le2\min\{1,\rho_{\mathrm B}\}
\le C\min\{1,\rho_{\mathrm B}\},
\]
using \(C\ge4\) and the nonnegativity of the capped rate.
% lean: CausalSmith.Stat.AnnotationRarearmFrontier.baseline_small_sample_risk
% lean: CausalSmith.Stat.AnnotationRarearmFrontier.baseline_small_sample_rate

\item
Finally, suppose \(n\ge6\). If \(\rho_{\mathrm B}\ge1\), the estimator and target both belong to \([-1,1]\), so
\[
\bigl(T^{\mathrm B}(\mathcal D)-\tau(P)\bigr)^2\le4
\]
for every realization. Hence
\[
\mathcal R_{\mathrm B}(P)
\le4
\le C
=
C\min\{1,\rho_{\mathrm B}\}.
\]
If \(\rho_{\mathrm B}<1\), the uncapped bound established above gives
\[
\mathcal R_{\mathrm B}(P)
\le C_{\mathrm{large}}\rho_{\mathrm B}
\le C\rho_{\mathrm B}
=
C\min\{1,\rho_{\mathrm B}\}.
\]
These cases cover every admissible experiment and every
\(P\in\mathcal M_{d,\epsilon}\), with the same positive universal constant \(C\).
% lean: CausalSmith.Stat.AnnotationRarearmFrontier.baseline_risk_le_four
% lean: CausalSmith.Stat.AnnotationRarearmFrontier.uniform_baseline
\end{enumerate}
\end{proof}
